\section{Теорема о размерности суммы подпространств}

\begin{theorem}[Формула Грассмана]
Для любых конечномерных подпространств $U_1$ и $U_2$ линейного пространства $V$ справедливо следующее соотношение для их размерностей:
\[
\dim(U_1 + U_2) = \dim U_1 + \dim U_2 - \dim(U_1 \cap U_2).
\]
\end{theorem}

\begin{proof}
Пусть $\dim(U_1 \cap U_2) = m$. Выберем базис пересечения: $b_1, b_2, \dots, b_m$.

\begin{enumerate}
    \item По теореме о неполном базисе, дополним эту систему до базиса в подпространстве $U_1$. Обозначим этот базис как $a_1, \dots, a_k, b_1, \dots, b_m$. Тогда $\dim U_1 = k + m$.
    \item Аналогично, дополним базис пересечения до базиса в подпространстве $U_2$. Обозначим его как $c_1, \dots, c_l, b_1, \dots, b_m$. Тогда $\dim U_2 = l + m$.
\end{enumerate}

Рассмотрим объединенную систему векторов:
\[
E = \{ a_1, \dots, a_k, b_1, \dots, b_m, c_1, \dots, c_l \}.
\]
Докажем, что $E$ является базисом суммы $U_1 + U_2$.

\begin{itemize}
    \item \textbf{Порождение:} Любой вектор $x \in U_1 + U_2$ представим в виде $x = x_1 + x_2$, где $x_1 \in U_1$, $x_2 \in U_2$. Так как $x_1$ линейно выражается через $\{a_i, b_j\}$, а $x_2$ --- через $\{c_p, b_j\}$, их сумма $x$ линейно выражается через всю систему $E$.
    
    \item \textbf{Линейная независимость:} Составим линейную комбинацию векторов системы $E$, равную нулю:
    \[
    \sum_{i=1}^k \alpha_i a_i + \sum_{j=1}^m \beta_j b_j + \sum_{p=1}^l \gamma_p c_p = 0. \quad (*)
    \]
    Перенесем слагаемые с векторами $c_p$ в правую часть:
    \[
    \sum_{i=1}^k \alpha_i a_i + \sum_{j=1}^m \beta_j b_j = - \sum_{p=1}^l \gamma_p c_p.
    \]
    Левая часть этого равенства принадлежит $U_1$, а правая --- $U_2$. Следовательно, этот вектор принадлежит пересечению $U_1 \cap U_2$. Значит, он может быть линейно выражен только через базис пересечения $b_1, \dots, b_m$:
    \[
    - \sum_{p=1}^l \gamma_p c_p = \sum_{j=1}^m \delta_j b_j \quad \Rightarrow \quad \sum_{j=1}^m \delta_j b_j + \sum_{p=1}^l \gamma_p c_p = 0.
    \]
    Так как система $\{b_1, \dots, b_m, c_1, \dots, c_l\}$ является базисом $U_2$, она линейно независима. Следовательно, все коэффициенты в этой комбинации равны нулю: $\gamma_p = 0$ для всех $p$ и $\delta_j = 0$ для всех $j$.
    
    Подставляя $\gamma_p = 0$ в исходное равенство $(*)$, получаем:
    \[
    \sum_{i=1}^k \alpha_i a_i + \sum_{j=1}^m \beta_j b_j = 0.
    \]
    В силу линейной независимости базиса $U_1$ ($\{a_1, \dots, a_k, b_1, \dots, b_m\}$), все $\alpha_i = 0$ и $\beta_j = 0$.
\end{itemize}

Таким образом, система $E$ линейно независима и порождает $U_1 + U_2$, то есть является его базисом. 
Подсчитаем количество векторов в базисе $E$:
\[
\dim(U_1 + U_2) = k + m + l.
\]
Заметим, что:
\[
\dim U_1 + \dim U_2 - \dim(U_1 \cap U_2) = (k + m) + (l + m) - m = k + m + l.
\]
Равенство доказано.
\hfill$\blacksquare$
\end{proof}
